% Lösungen Einheit 1 – Prozentsatz berechnen
\einheitenkopf*{Einheit 1 von 4 · Prozentsatz berechnen}
\erg{11}{a) 30 Kinder \quad b) 10 Aufgaben \quad c) 40\,€ \quad d) 1 Liter \quad e) alle Murmeln (nicht die blauen)}
\erg{12}{a) 1 Strich bei 50\,\% \quad b) 3 Striche bei 25, 50, 75\,\% \quad c) 9 Striche, alle 10\,\% \quad d) 4 Striche bei 20, 40, 60, 80\,\%}
\erg{13}{Beispiel 40\,\% \quad a) 1 Kästchen $=$ 10\,\%, 70\,\% \quad b) 1 Kästchen $=$ 2\,€ $=$ 10\,\%, 30\,\% \quad c) 1 Kästchen $=$ 1\,h $=$ 25\,\%, 25\,\% \quad d) 1 Kästchen $=$ 5\,m $=$ 10\,\%, 90\,\%}
\erg{14}{a) 10\,\% \quad b) 30\,\% \quad c) 50\,\% \quad d) 70\,\% \quad e) Das Ganze (100) bleibt gleich; der Teil wächst jedes Mal um 20, der Prozentsatz auch.}
\erg{15}{a) $\frac{13}{25} = \frac{52}{100} = 52\,\%$ \quad b) $27 : 60 = 0{,}45 = 45\,\%$ \quad c) $39 : 47 \approx 0{,}8298 \approx 83{,}0\,\%$}
\erg{16}{z.\,B. a) 1 von 2 \quad b) 1 von 4 \quad c) 1 von 10 \quad d) 3 von 5 \quad (richtig ist jedes Paar mit Teil : Ganzes $=$ Prozentsatz)}
\erg{17}{a) Ganzes $18 + 12 = 30$; $18 : 30 = 0{,}6 = 60\,\%$ \quad b) Rabatt $50 - 40 = 10$\,€; $10 : 50 = 0{,}2 = 20\,\%$ \quad c) Ganzes $5 + 19 = 24$; $5 : 24 \approx 0{,}2083 \approx 20{,}8\,\%$}
\erg{18}{a) Tim teilt das Ganze durch den Teil; richtig $9 : 20 = 0{,}45 = 45\,\%$ \quad b) $7 : 25 = 0{,}28 = 28\,\%$}
\erg{19}{a) Mia teilt durch die anderen Bonbons statt durch alle 20; richtig $4 : 20 = 0{,}2 = 20\,\%$ \quad b) Ganzes $9 + 21 = 30$; $9 : 30 = 0{,}3 = 30\,\%$}
\erg{20}{a) $6 : 30 = 0{,}2 = 20\,\%$: Der Teil wird durch das Ganze geteilt; $30 : 6 = 5 = 500\,\%$ wäre mehr als das Ganze. \quad b) Stimmt: Ein Teil ist höchstens so groß wie sein Ganzes, der Anteil also höchstens $1 = 100\,\%$.}
\erg{21}{a) $18 : 24 = 0{,}75 = 75\,\%$ \quad b) $21 : 30 = 0{,}7 = 70\,\%$ \quad c) Lea trifft besser (75\,\% > 70\,\%); Ben hat nicht recht, er hat nur öfter geworfen.}
